\section{Iterative DPLL}

Now that we've implemented DPLL to the best of our ability, it's time to 
take it one step further. We will now implement it iteratively rather than 
recursively. Why? Remember how I said in the beginning that DPLL is the foundation 
for modern SAT solvers?
\footnote{I can't remember if I actually said that}
Well iterative DPLL is the foundation for the next chapter's algorithm: CDCL.
Implementing iterative DPLL is imperative in order to create a modern SAT solver
worth its salt.
\footnote{Honestly, ``Salt'' is a better name than ``Saturn.'' I might change it.}

As said in the beginning of the chapter, DPLL forces us to traverse the decision 
tree as well, a tree. And that is still true, even in iterative DPLL. What changes 
now is that instead of an implicit path, we now have an explicit one. Suppose we 
are on the first decision level, so we decide variable $1 = T$, and that forces 
$3 = F$, $5 = F$, and $6 = T$. In recursive DPLL, we assign these variables into 
the \textit{assignment} or \textit{vars} field, but we don't really use them until 
we get to the decision level that matches their index or the evaluation phase. 

With iterative DPLL, there are still decision levels, but the representation is 
now very, very different bcause instead of a decision \textit{tree}, we now use a 
decision \textit{trail}. The trail is reminiscent of a queue in the sense that 
we only operate at the end of the data structure. It contains some information 
that we need in order to continue mimicking the behavior of recursive DPLL. Each 
trail entry contains the literal that was forced, the decision level it's a part of,
and the clause that forced that decision. We append these trail entries every single 
time we force a variable's assignment. The next question is how do we know which 
variables were decided for and which variables were inferred from propagation?
Well, for decisions, there is no reason clause because there was no reason to 
begin with other than ``because I said so.'' Therefore, the variables that were 
decided for don't get a reason clause, but the ones that were inferred do. This 
is great because it makes it easy to tell where the beginning of one decision 
starts and where it ends.

However, it would be nice to have some way to keep track of where the decisions
are in the trail other than scanning for entries without a reason clause. Because 
of that, we will also keep another trail specifically for the decision markers.
This is also helpful because then we can store whether or not we've ``tried the 
other branch'' when we backtrack. 

While on the topic of backtracking, it would be nice to discuss it some more as well. 
Since we can't simply go up the call stack and undo variable assignments, we have to 
do it in the iterative way. When we backtrack, we first backtrack up to the decision
entry. This is done by popping off the trail and undoing the propagated assignments. 
Once that's done, you then check if you've ``tried the other branch'', and if you 
haven't, you flip the decision literal's assignment and then retry. If you have 
already tried the other branch, you pop off the decision entry, and then repeat 
this process to the next decision entry. This process continues until you either 
find a decision that has not been switched yet, or you run out of decision 
entries.

With most of the exposition out of the way, I think it would be good to 
look at the changes we will make to our interface and algorithms, starting 
with our SAT solver.

\begin{lstlisting}
    public interface:
        SATSolver(string input)
        bool solveCNF()

    private interface:
        enum var
        {
            False
            True
            Unassigned
        }

        TrailEntry {
            int lit,
            int clauseIdx
        }

        DecisionLevel {
            int idx,
            bool triedOtherBranch
        }

        void parse(string equation)
        bool propagate()
        bool createWatched()

        int numVars
        int numClauses
        int trailHead

        int[][] clauses
        bool[] satisfied
        
        TrailEntry[] trail
        DecisionLevel[] trailStarts

        optional var[] vars
        map<int, int[]> var_to_clause
        map<int, (int, int)> clause_to_vars
\end{lstlisting}

The first notable change is that \textit{propagate()} now returns a boolean instead 
of a delta. This is because it will no longer be responsible for undoing propagation 
in response to a contradiction, thus achieving zero-cost backtracking. Next, 
we now have a \textit{trail} and a \textit{trailHead}, which points to the last 
made decision entry. \textit{trailStarts} keeps track of all the indices of the 
decision entries. One last thing to note is that the \textit{TrailEntry} data 
structure uses indices for clauses instead of the actual clause itself. 
There's a very good reason for this, and I will reveal the answer in the next chapter,
but until then, why do you think it might be that way? Hint: it's not because 
indices are easy for accessing. 

The next change we need to make is to our algorithm to create watched literals. 
It's the same for basically $90\%$ of the algorithm, except for the unit clause 
case, which I'll highlight.

\begin{algorithm}[H]
\caption{Modified createWatched}\label{alg:21}
\begin{algorithmic}[1]
\Procedure{createWatched}{}

\For{each $Clause$ in $Clauses$}
    \If{$Clause$ has at least 2 literals}
        \State ...
    \ElsIf{$Clause$ is has 1 literal}
        \State lit $\gets$ Clause[0]
        \State ...
        \State cIdx $\gets$ $\Call{indexOf}{Clauses, Clause}$

        \State $\Call{append}{trail, (lit, cIdx)}$
        \State ...
    \Else
        \State return $false$
    \EndIf
\EndFor
\State return $true$

\EndProcedure
\end{algorithmic}
\end{algorithm}

Rather than enqueuing the variable, we create a new trail entry for each unit 
clause. These ones are all part of the \textbf{root} decision level, which we 
mentally reserve for the 0th decision level. If we backtrack up to this level (if it exists)
or before it, we've encountered a root level contradiction, so the equation is 
unsatisfiable. Because not every equation has unit clauses, it's important to 
remember that this decision level may not always exist, and so we should be 
keep that in mind for our solving algorithm. 

The next algorithm we're going to change is the \textit{propagate()} algorithm 
because the changes are also very minimal.

\begin{algorithm}[H]
\caption{Modified propagate}\label{alg:22}
\begin{algorithmic}[1]
\Procedure{propagate}{assignment[1..n]}

\While{trailHead < $\Call{size}{trail}$}
    \State prop $\gets$ trail[trailHead].lit
    \State trailHead $\gets$ trailHead + 1
    \State ...

    \State watchedClauses $\gets$ varToClause[prop]

    \For{each $cIdx$ in watchedClauses}
        \State Clause[1..t] $\gets$ Clauses[cIdx]
        \State watches $\gets$ clauseToVar[cIdx] \Comment pair of clause indices
        \State ...

        \If{relocated = False}
            \State ...
            \State otherWatch $\gets$ Clause[secondWatch]
            \State oIdx $\gets \Call{abs}{otherWatch}$

            \If{assignment[oIdx] = Unassigned}
                \State ...
                \State $\Call{append}{trail, (otherWatch, cIdx)}$
            \Else
                \If{\Call{evalsFalse}{otherWatch} = false}
                    \State continue onto the next loop iteration
                \Else 
                    \State $\Call{undoPropagation}{delta}$
                    \State $\Call{clear}{unitProp}$
                    \State return $false$
                \EndIf
            \EndIf
        \EndIf
    \EndFor
\EndWhile

\State return $true$

\EndProcedure
\end{algorithmic}
\end{algorithm}

The main difference is that we just \textit{trailHead} to iterate through the trail, and 
instead of enqueuing, we just append to the trail. We also return True/False depending on 
whether or not propagation succeeded.

Lastly, we look at our \textit{solve()} method, which has the most notable changes.
There are two ways to implement this, depending on what makes sense to you. I call 
them ``assignment before propagation'' and ``propagation before assignment.'' I
call them these because that's what they are. Before showing the pseudocode, I 
want to talk about the algorithm a bit more. The beginning is mostly the same:
create our watched literals, check that it succeeded, return False if it didn't. 
Otherwise, we do a root level propagation and then check if it succeeds as well.
However, the main confusion comes from handling the case of when there is a 
0th decision level and when there isn't. There are multiple ways to go about this.
After the beginning, we enter the solving loop. 

There are two main parts: assigning and backtracking, with a call to the propagate 
method separating them. Assigning is now scanning for the first Unassigned variable,
forcing it True, creating a new trail entry, and then updating the trail head. I'm 
sure there's a better way to keep track of the next unassigned variable other than 
scanning the entire assignment list on each loop iteration, however, for simplicity's 
sake, we're implementing it the naive way. 

Since we've already discussed backtracking, we will now get into the algorithm 
itself:

\begin{algorithm}[H]
\caption{Modified solving loop}\label{alg:23}
\begin{algorithmic}[1]
\Procedure{solve}{}

\State setupFailed $\gets \Call{setup}{}$

\If{setupFailed = true}
    return $false$
\EndIf

\While{$true$}
    \State succeeded $\gets \Call{propagate}{assignment}$
    \If{succeeded = $false$}
        \State \Call{backtrack}{}

        \If{\Call{size}{trailStarts} = 0}
            \State return $false$
        \EndIf

        \State continue
    \EndIf

    \If{\Call{assign}{assignment} = true}
        \State return $true$
    \EndIf

\EndWhile

\EndProcedure
\end{algorithmic}
\end{algorithm}

\begin{algorithm}[H]
\caption{Setup Solve}\label{alg:24}
\begin{algorithmic}[1]
\Procedure{setup}{}

\State assignment[1..numVars]
\For{each $var$ in $assignment$}
    \State $var \gets$ Unassigned
\EndFor

\If{\Call{createWatched}{assignment} = false}
    \State return $true$
\EndIf

\State \Call{append}{trailStarts, (0, true)}

\If{\Call{size}{trail} > 0}
    \If{\Call{propagate}{assignment, 0} = $false$}
        \State return $true$
    \EndIf
\EndIf

\State return $false$

\EndProcedure
\end{algorithmic}
\end{algorithm}

\begin{algorithm}[H]
\caption{Backtracking code}\label{alg:25}
\begin{algorithmic}[1]
\Procedure{backtrack}{}

\While{\Call{size}{trailStarts} > 0}
    \State decision $\gets$ most recent decision level

    \While{\Call{size}{trail} > decision.idx + 1}
        \State lit $\gets$ last trail entry's literal
        \State assignment[\Call{abs}{lit}] $\gets$ Unassigned 
        \State erase last trail entry
    \EndWhile

    \If{decision.TriedFalse = false}
        \State entry $\gets$ last trail entry
        \State entry.lit $\gets$ -entry.lit
        \State assignment[\Call{abs}{lit}] $\gets$ $false$
        \State decision.TriedFalse $\gets$ $true$
        \State trailHead $\gets$ decision.idx
        \State break
    \Else
        \State lit $\gets$ last trail entry's lit
        \State assignment[\Call{abs}{lit}] $\gets$ Unassigned 
        \State erase last trail entry
        \State erase last trailStarts entry
    \EndIf
\EndWhile

\EndProcedure
\end{algorithmic}
\end{algorithm}

\begin{algorithm}[H]
\caption{Assignment code}\label{alg:26}
\begin{algorithmic}[1]
\Procedure{assign}{}

\State firstUnassigned $\gets$ -1
\For{$i \gets 1$ up to numVars}
    \If{assignment[i] = Unassigned}
        \State firstUnassigned $\gets$ i
        \State break
    \EndIf
\EndFor

\If{firstUnassigned = -1}
    \State return $true$
\EndIf

\State assignment[firstUnassigned] = $true$
\State \Call{append}{trail, (firstUnassigned, -1)}
\State \Call{append}{trailStarts, (trail size - 1, false)}
\State trailHead $\gets$ trailStarts.last.idx

\State return $false$

\EndProcedure
\end{algorithmic}
\end{algorithm}

Here is the pseudocode for solving, written as ``propagate then assign." Quick question:
when there's no 0th decision level and we propagate, how do we know that none of the 
propagated variables will have the wrong decision level? Ie, how do we know that they'll 
have a decision level of $0$ instead of $1$, since we haven't made any decisions yet?
This is kind of a trick question, the answer is that if there is no 0th level,
\textit{propagate()} will do nothing because the while-loop condition will immediately 
return False, and so the method will terminate, assigning nothing! 

Instead of showing you the code for ``assign then propagate'' (it is NOT just taking 
the assigning code and moving it up before the call to \textit{propagate()}), let's 
run through how this works. When there's a 0th level, we obviously want to mark it 
in our \textit{trailStart}, but why do we want to mark that the ``other branch'' has 
already been tried? Well that's because those assigned by the root level propagation 
are forced into those values, so it doesn't make sense to swap their values to be 
their opposites since it will immediately be an unsatisfiable assignment anyways. 

Next, in the main solving loop, the backtracking code is just ``grab the last 
decision index, undo until the root decision, if it's already tried False, undo 
the decision entirely and repeat with the previous one, and if not, try False 
and continue.''

Lastly, let's talk about doing assign then propagation. Here's why it doesn't
make total sense to start out with propagation:

\begin{itemize}
    \item we do a root level propagation if there needs to be one.
    \item if there is no root level propagation to be done, it still
            doesn't make sense to immediately propagate since there's 
            nothing to propagate.
    \item if there was a root level propagation, it doesn't make sense 
            to do another one immediately since we've just completed one.
\end{itemize}

So if you wanted to do assign and then propagate, how would that be done? 
Earlier, I said that it's not just taking the code for assigning and then 
moving it before propagation, but why? Well, suppose we fail a propagation. 
We backtrack, and then suppose we haven't tried the other branch for that 
decision. We encounter a \textit{continue} statement and we go back to the 
top. If we simply move the assigning code up to the top of the loop,
what's the issue? Well remember, we just flipped a decision, so we should
actually be \textbf{\textit{propagating}} instead of assigning. If we 
assign immediately after flipping the decision, we now have two decisions 
at the same time, and that's bad because what if they are contradictory to 
one another, ie one would force a variable False and the other True? 
Unfortunately, the propagation code wouldn't catch that because it assumes 
everything in the trail is valid, but we just made the trail invalid. So 
to the propagation code, instead of it catching that a variable is forced to be 
False, but it's True, or vice versa, it may just ignore it because it focuses 
mainly on Unassigned and False literals.

So if we want to do assign then propagate, how do we make sure that 
we don't accidentally push something onto the trail when we shouldn't have?
We need to wrap the assigning code in a conditional statement that checks that 
\textit{trailHead == trail.size}, if it doesn't, then that means that there's 
something to be propagated on because we're not at the end of the trail, but 
if it does, then we've reached the end, so we should check if we've assigned 
everything and can be done, or if we need to assign something new.

\subsection{Iterative DPLL - Results}
Here are the results after this huge refactor:

(Assignment before propagation)\\
uf20.91: \textbf{328ms}, down 21ms from \textbf{349ms} by Watched Literals\\
UUF50.218.1000: \textbf{6369ms}, down 761ms from \textbf{7130ms} by Watched Literals

(Propagation before assignment)\\
uf20.91: \textbf{324ms}\\
UUF50.218.100: \textbf{6365ms}

\begin{FVerbatim}
On all 1000 equations in uf20-91:

Iterative DPLL (AbP):      328ms (~0.33 seconds).
Iterative DPLL (PbA):      324ms (~0.32 seconds).
Watched Literals:          349ms (~0.35 seconds).
Pure Literal Elimination:  1548ms (~1.55 seconds).
Unit Propagation:          868ms (~0.87 seconds).
Recursive Backtracking:    1,540,613ms (~25.68 minutes).
Brute Force:               1,770,851ms (~29.51 minutes).
Short-Circuited (Basic):   49,874ms (~50 seconds).
Parallel:                  10,589ms (~10.5 seconds).

On all 1000 equations in UUF50.218.1000:

Iterative DPLL (AbP):     6,369ms (~6.37 seconds).
Iterative DPLL (PbA):     6,365ms (~6.37 seconds).
Watched Literals:         7,130ms (~7.13 seconds).
Pure Literal Elimination: 231,107ms (~3.85 minutes).
Unit Propagation:         140,041ms (~2.34 minutes).
\end{FVerbatim}

Should you do propagation before assignment simply because it's faster? It 
really doesn't matter, you do you. However, for the rest of this book, we will be 
using assignment before propagation because it's easier for me to understand, and it 
follows online pseudocode examples better than propagate before assignment.

The speed gain on satisfiable problems is very miniscule, however, if you ever want 
proof that iterative algorithms are faster than recursive ones, the unsatisfiable 
problem time decrease is pretty good evidence.